% These rather complex beasties are the various macros used to make up
% The information on a product.

% ---------------------------------------- Defining a product
% This command sets up a section for a species.
%
\newcommand{\thumbmark}{}
\newcommand{\thumbbackground}{black}
%
\newcommand{\speciestitle}[5][]{%
\cleardoublepage\suppressfloats[t]
\renewcommand{\thumbmark}{#3}%
\addtolength{\thumbpos}{-0.45in}%
\section{#2}
\begin{description}
\item[Swath name:] \texttt{#4} #1
#5
\end{description}}
%
\newlength{\thumbpos}
\newcommand{\makethumbmark}[2]{%
  \begin{tikzpicture}[overlay]
    \node at (#1,\thumbpos) [rotate=#2,anchor=south]
    {\colorbox{black}{\parbox{1.5cm}{\centering%
        \color{white}\sffamily\bfseries\large\selectfont\rule[-0.5ex]{0pt}{2ex}\thumbmark}}};%
  \end{tikzpicture}}
%
\newcommand{\startthumbmarks}{%
  \setlength{\thumbpos}{9.2in}%
  \lfoot[\makethumbmark{-1cm}{90}\thepage]{\myshorttitle}%
  \rfoot[Aura Microwave Limb Sounder]{\thepage\makethumbmark{1cm}{-90}}%
} \newcommand{\killthumbmarks}{%

  \lfoot[\thepage\hspace{1cm}\tinyrcs]{\myshorttitle} % should be \@shortitle
  \rfoot[EOS Microwave Limb Sounder]{\tinyrcs\hspace{1cm}\thepage}}
\newcommand{\email}[1]{\textbf{Email:} $<$\texttt{#1}$>$}
%
% Change the figures/table counters
\makeatletter
\def\renewcounter#1{%
    \@ifundefined{c@#1}
    {\@latex@error{counter #1 undefined}\@ehc}%
    \relax
    \let\@ifdefinable\@rc@ifdefinable
    \@ifnextchar[{\@newctr{#1}}{}}
\makeatother
\renewcounter{figure}[section]
\renewcounter{table}[section]
\renewcommand{\thefigure}{\thesection.\arabic{figure}}
\renewcommand{\thetable}{\thesection.\arabic{table}}
%
\newcommand{\screeningrule}[1]{\textbf{#1}\par}

% -------------------------- Averaging kernels
\newcommand{\extracaption}{}

\newcommand{\avkcaption}[2]{%
  Typical two-dimensional (vertical and horizontal along-track)
  averaging kernels for the MLS \thisv #2#1 data at 70\degsym N (upper)
  and the equator (lower); variation in the averaging kernels is
  sufficiently small that these are representative of typical profiles.
  Colored lines show the averaging kernels as a function of MLS
  retrieval level, indicating the region of the atmosphere from which
  information is contributing to the measurements on the individual
  retrieval surfaces, which are denoted by plus signs in corresponding
  colors.  The dashed black line indicates the resolution, determined
  from the full width at half maximum (FWHM) of the averaging kernels,
  approximately scaled into kilometers (top axes).  (Left) Vertical
  averaging kernels (integrated in the horizontal dimension for five
  along-track profiles) and resolution.  The solid black line shows the
  integrated area under each kernel (horizontally and vertically);
  values near unity imply that the majority of information for that MLS
  data point has come from the measurements, whereas lower values imply
  substantial contributions from a priori information.  (Right)
  Horizontal averaging kernels (integrated in the vertical dimension)
  and resolution.  The horizontal averaging kernels are shown scaled
  such that a unit averaging kernel amplitude is equivalent to a factor
  of 10 change in pressure.}

\newcommand{\vertavkcaption}[2]{%
  Typical vertical averaging kernels for the MLS \thisv #2#1 data at
  70\degsym N (left) and the equator (right); variation in the averaging
  kernels is sufficiently small that these are representative of typical
  profiles.  Colored lines show the averaging kernels as a function of
  MLS retrieval level, indicating the region of the atmosphere from
  which information is contributing to the measurements on the
  individual retrieval surfaces, which are denoted by plus signs in
  corresponding colors.  The dashed black line indicates the vertical
  resolution, determined from the full width at half maximum (FWHM) of
  the averaging kernels, approximately scaled into kilometers (top
  axes).  The solid black line shows the integrated area under each
  kernel; values near unity imply that the majority of information for
  that MLS data point has come from the measurements, whereas lower
  values imply substantial contributions from a priori information.  The
  low signal to noise for this product necessitates the use of
  significant averaging (e.g., monthly zonal mean), making horizontal
  averaging kernels largely irrelevant.}

\newcommand{\onestandardavk}[2]{%
\begin{figure}%
\begin{center}%
\includegraphics[width=0.8\linewidth]{figures/avk-#1}%
\end{center}
\caption{\avkcaption{#2}{}\extracaption}
\label{fig:Avk#1}%
\end{figure}}

\newcommand{\oneverticalavk}[2]{%
\begin{figure}%
\begin{center}%
\includegraphics[width=0.8\linewidth]{figures/avk-#1}%
\end{center}
\caption{\vertavkcaption{#2}{}\extracaption}
\label{fig:Avk#1}%
\end{figure}}

% ------------------------------------ Other boilerplate
\newcommand{\standardrangestatement}{Values outside this range are not
  recommended for scientific use.}
\newcommand{\standardprecisionrule}{%
\item[Estimated precision:] \screeningrule{Only use values for which the
    estimated precision is a positive number.} Values where the \emph{a
    priori} information has a strong influence are flagged with negative
  precision, and should not be used in scientific analyses (see
  Section~\ref{sec:NegativePrecision}).}
\newcommand{\evenstatusrule}{%
\item[Status flag:] \screeningrule{Only use profiles for which the
    `Status' field is an even number.}  Odd values of \texttt{Status}
  indicate that the profile should not be used in scientific studies.
  See Section~\ref{sec:StatusQuality} for more information on the
  interpretation of the \texttt{Status} field.}
\newcommand{\cloudsokrule}{\screeningrule{%
    Profiles identified as being affected by clouds can be used with no
    restriction.}}
\newcommand{\qualityrule}[1]{\screeningrule{Only
    profiles whose `Quality' field is greater than #1 should be used.}}
\newcommand{\convergencerule}[1]{\screeningrule{Only profiles whose
    `Convergence' field is less than #1 should be used.}}

% ----------------------------------- Table flipflopping
\colorlet{tableoddbackground}{black!10}
\colorlet{tableevenbackground}{white}
\newif\ifoddrow\oddrowtrue
\newcommand{\flipfloprow}{%
  \ifoddrow
  \rowcolor{tableoddbackground}
  \global\oddrowfalse
  \else
  \rowcolor{tableevenbackground}
  \global\oddrowtrue
  \fi}
\newcommand{\samecolorrow}{%
  \ifoddrow
  \rowcolor{tableevenbackground}
  \else
  \rowcolor{tableoddbackground}
  \fi}  

%%% Local Variables: 
%%% mode: latex
%%% TeX-master: "v4-x-report"
%%% End: 
